\subsection{固定百分位参照的数学性质与评分解释}
中秩百分位变换同时解决了量纲不同和极端值影响的问题，但也引入了需要说明的性质。对某个方向已统一的字段，将 A1 中相等的观测分为若干组，第 $v$ 组有 $m_v$ 条，前面累计有 $a_v$ 条，则该组每条记录的百分位为 $(a_v+m_v/2)/n$。对全部记录求和可得
\begin{equation}
\frac{1}{n}\sum_{i=1}^{n}r_{ij}
=\frac{1}{n^2}\sum_v m_v\left(a_v+\frac{m_v}{2}\right)
=\frac12.
\label{eq:q1-rank-mean}
\end{equation}
最后一个等式来自有序观测对计数：不同组的配对贡献与同组的一半贡献合起来为 $n^2/2$。因此，只要所有记录使用同一参考集、权重非负且和为 $1$，A1 样本总体的平均 $Q$ 必然为 $0.5$。该性质与数据实际是否适合训练没有一一对应关系；平均数接近 $0.5$ 不能用来验证评分“准确”。对 A2/A3 使用 A1 参照时，上述恒等式不再强制成立，扩展集平均分才可以与 A1 同名域作同标度对照。

固定参照还能避免重新排名产生的伪稳定性。若每个新增集合都在内部重新转成百分位，其总体平均分仍会回到 $0.5$，即使新集合的原始质量信号整体下降，均分也不一定反映下降。本文保存 A1 的中位数、截断阈值和经验分布，使后来的样本始终回答同一个问题：相对于 A1，该样本处在何种位置。超出 A1 支撑范围的数值被截在边界，因此边界附近的评分只表明超出参照，不提供超出程度的线性测量。

该变换对严格单调递增重参数化保持排序不变。例如在不发生截断时，正向字段先取对数再计算经验百分位，不会改变记录次序；长度字段的实际作用变化主要来自第 $95$ 百分位截断形成的并列值。反向字段则通过方向转换明确“越低越好”。这些处理应由字段含义决定，不能因为某种方向让结果更符合预期就临时调整。缺失填补同样只是保证评分可计算的规则，并不表示缺失记录拥有真实的中等质量。

\subsection{组间权重敏感性与局部冲突的区别}
记三组平均百分位分别为 $M_i,S_i,T_i$，则综合分可以写为
\begin{equation}
Q_i(\bm a)=a_M M_i+a_S S_i+a_T T_i,
\qquad a_M+a_S+a_T=1,\quad a_g\geq0.
\label{eq:q1-group-score}
\end{equation}
其中 $S_i$ 为 DSIR 组、$T_i$ 为结构组。对于两个域 $d,e$，域均值之差等于三组域均值差与权重的内积。若一个域在三个组上都不低于另一域，则所有非负权重下的域序都一致；若各组的优势方向不同，改变权重就可能改变域序。因此，域排序稳定性反映了指标之间是否形成一致判断，不能仅由主方案某两个得分的差距判断。

图~\ref{fig:q1-weight-robustness} 展示三组权重各限于 $0.2$--$0.5$、步长为 $0.05$ 的 $36$ 个可行设置。每个点表示一个事先定义的权重组合，并非一次随机抽样。所得 $Q_0$ 介于 $0.45514$ 与 $0.51139$；C4 在 $28$ 个设置中排名最高，arxiv 在 $8$ 个设置中最高。书籍在 $31$ 个设置中最低，github 在 $5$ 个设置中最低。该有限网格刻画一组偏好变化下的结果范围，不是某域排名的概率或置信区间。

\begin{figure}[!htbp]
\centering
\includegraphics[width=0.88\textwidth]{fig_q1_weight_robustness.pdf}
\caption{组间权重变化下的基线质量与域序稳定性；每点为一个确定性设置}
\label{fig:q1-weight-robustness}
\end{figure}

组内 CRITIC 对照的域序完全不变，而组间权重改变会出现明显重排，说明进一步校准应优先回答“模型评分、DSIR 与结构信号各应占多大比例”。单纯增加赋权算法复杂度，无法替代与下游训练目标相对应的标签。对于当前题目，三组平衡主方案的优势在于规则透明、不会因某组指标数量较多而自动获得更高总权重；它仍然是一个可以被新数据修正的选择。

局部冲突与综合质量描述不同现象。表~\ref{tab:q1-conflict-domains} 同时列出发生数与域内比例：C4 不仅综合质量较高，其冲突率也最高。高综合分不会排除少数记录同时具有广告和教育价值，二者并不矛盾。书籍只有 $171$ 条，一条冲突就对应约 $0.585\%$；相比之下，github 的一条冲突对应 $0.010\%$。仅按冲突数比较，或忽略分母直接排列比率，都可能误判数据治理的优先级。

\begin{table}[!htbp]
\centering
\songti\zihao{-4}
\setlength{\tabcolsep}{4pt}
\caption{A1 各域冲突计数与域内发生率}
\label{tab:q1-conflict-domains}
\begin{tabular}{lrrr}
\toprule
来源域 & 样本数 & 冲突数 & 域内冲突率 \\
\midrule
c4 & 10000 & 103 & 1.0300\% \\
book & 171 & 1 & 0.5848\% \\
commoncrawl & 9640 & 48 & 0.4979\% \\
wikipedia & 10000 & 2 & 0.0200\% \\
github & 10000 & 1 & 0.0100\% \\
arxiv & 1419 & 0 & 0.0000\% \\
stackexchange & 10000 & 0 & 0.0000\% \\
\bottomrule
\end{tabular}
\end{table}


实际使用时可保留两个分开的信息：以 $Q$ 表示整体相对可用性，以冲突标记触发复核或专项处理。当前 $\delta$ 降权只用于敏感性分析，尚未计入训练质量收益。若后续获得清洗前后训练实验，应在相同 token 预算、相同模型结构与训练配方下比较，而不是仅以降权后的评分下降或排名稳定作为治理有效的证据。
